\paragraph{\texttt{auger1989segment}: \citet{auger1989segment}.}
Auger and Lawrence develop dynamic-programming algorithms for identifying segment neighborhoods under broad segment-fit criteria, including least squares and maximum likelihood. The paper is part of the classical algorithmic lineage for Bellman-style segmentation and shows that an optimal-partition max-sum recurrence is not itself a BayesBreak contribution. BayesBreak differs by using integrated segment likelihoods in Bayesian recursions for normalization and posterior marginals, while retaining a separate max-sum recursion for the joint MAP partition. The paper is used as algorithmic background rather than as a source for Bayesian posterior calculations. \textit{Verification status: title, venue, volume, pages, and DOI independently checked against the publisher record in Phase~6.}
